\section*{Chapter \ref{ch:b3:group-theory-applied} --- Group Theory Applied}

\begin{solution}{exo:b3:group-theory-applied:1}
$n_{A_1} = \frac14(5 - 1 + 3 + 1) = 2$, $n_{A_2} = \frac14(5 - 1 - 3 - 1) = 0$,
$n_{B_1} = \frac14(5 + 1 + 3 - 1) = 2$, $n_{B_2} = \frac14(5 + 1 - 3 + 1) = 1$:
$\Gamma = 2A_1 + 2B_1 + B_2$, of dimension 5.
\end{solution}

\begin{solution}{exo:b3:group-theory-applied:2}
Same atoms-in-place count as water: $\Gamma_{\mathrm{vib}} = 2A_1 + B_1$. $A_1$ ($z$) and
$B_1$ ($x$) are \omterm{def:b3:group-theory-applied:activity}{IR active}, and both are also \omterm{def:b3:group-theory-applied:activity}{Raman active} ($x^2$, $xz$): three
bands in each spectrum.
\end{solution}

\begin{solution}{exo:b3:group-theory-applied:3}
$B_1\otimes B_2$: \omterm{def:b3:point-groups:representation}{characters} $(1, 1, -1, -1)$ $= A_2$. $A_2\otimes B_1$: $(1, -1, -1,
1) = B_2$. In $C_{3v}$, $E\otimes E$ has \omterm{def:b3:point-groups:representation}{characters} $(4, 1, 0)$; reduction: $A_1 + A_2 +
E$.
\end{solution}

\begin{solution}{exo:b3:group-theory-applied:4}
Infrared: the two $t_2$ modes, 3019 and \qty{1306}{cm^{-1}}. Raman: all four ($a_1$,
$e$, $t_2$ all transform as quadratic functions). No \omterm{def:b3:point-groups:inversion}{centre of inversion}, so
coincidences are allowed.
\end{solution}

\begin{solution}{exo:b3:group-theory-applied:5}
Atoms in place $(4, 1, 2, 4, 1, 2)$ times $\chi_{xyz} = (3, 0, -1, 1, -2, 1)$ give
$\chi_{3N} = (12, 0, -2, 4, -2, 2)$, which reduces to $A_1' + A_2' + 3E' + 2A_2'' +
E''$. Remove $E' + A_2''$ (translations) and $A_2' + E''$ (rotations): $A_1' + 2E'
+ A_2''$. IR: $2E' + A_2''$, three bands; Raman: $A_1' + 2E'$, three bands. $A_2''$
(out-of-plane bend) is IR only, $A_1'$ (symmetric stretch) Raman only, the $E'$
modes appear in both.
\end{solution}

\begin{solution}{exo:b3:group-theory-applied:6}
IR: the two $T_{1u}$ modes, two bands. Raman: $A_{1g}$, $E_g$, $T_{2g}$, three bands.
$T_{2u}$ is silent. \ce{SF6} has a \omterm{def:b3:point-groups:inversion}{centre of inversion}: mutual exclusion.
\end{solution}

\begin{solution}{exo:b3:group-theory-applied:7}
Under $E$, $C_3$, $C_3^2$ the function $h_2$ goes to $h_2$, $h_3$, $h_1$ (\omterm{def:b3:point-groups:representation}{characters} 2,
$-1$, $-1$); reflections contribute 0. $\hat P_Eh_2 \propto 2h_2 - h_3 - h_1$.
Its overlap with $2h_1 - h_2 - h_3$ (atomic overlaps neglected) is $-3$, and the
latter has norm $\sqrt6$; subtracting the projection, $(2h_2 - h_1 - h_3) +
\frac12(2h_1 - h_2 - h_3) = \frac32(h_2 - h_3)$.
\end{solution}

\begin{solution}{exo:b3:group-theory-applied:8}
fac ($C_{3v}$, \omterm{def:b3:point-groups:representation}{characters} $3, 0, 1$): $A_1 + E$, two IR bands. mer ($C_{2v}$,
\omterm{def:b3:point-groups:representation}{characters} $3, 1, 3, 1$): $2A_1 + B_1$, three bands. \ce{M(CO)6} ($O_h$): $A_{1g} + E_g
+ T_{1u}$, one IR band ($T_{1u}$).
\end{solution}

\begin{solution}{exo:b3:group-theory-applied:9}
\omterm{def:b3:point-groups:representation}{Characters} $6, 0, 0, 2, 2, 0, 0, 0, 4, 2$: only $E$, the $C_4$ and $C_2$ along the
axes (two ligands left in place), the three $\sigma_h$ (four) and the six $\sigma_d$
(two) keep ligands in place. The reduction gives $A_{1g} + E_g + T_{1u}$. Metal $s$
($a_{1g}$), $d_{z^2}$ and $d_{x^2-y^2}$ ($e_g$), and $p_x$, $p_y$, $p_z$ ($t_{1u}$) combine;
$t_{2g}$ stays non-bonding.
\end{solution}

\begin{solution}{exo:b3:group-theory-applied:10}
$n_i = \frac1h(1\cdot h\cdot\chi_i(E)) = d_i$. The dimension of the regular
representation is $h$, and also $\sum_in_id_i = \sum_id_i^2$.
\end{solution}

\begin{solution}{exo:b3:group-theory-applied:11}
In $D_{2h}$ (molecule along $z$), atoms in place $(4, 4, 0, 0, 0, 0, 4, 4)$, $\chi_{3N}
= (12, -4, 0, 0, 0, 0, 4, 4)$, reducing to $2A_g + 2B_{2g} + 2B_{3g} + 2B_{1u} + 2B_{2u}
+ 2B_{3u}$. Remove translations $B_{1u} + B_{2u} + B_{3u}$ and the two rotations
$B_{2g} + B_{3g}$: $2A_g + B_{2g} + B_{3g} + B_{1u} + B_{2u} + B_{3u}$, seven modes.
In $D_{\infty h}$: $2\Sigma_g^+ + \Pi_g + \Sigma_u^+ + \Pi_u$. IR: $\Sigma_u^+$ and $\Pi_u$
(two bands); Raman: $2\Sigma_g^+$ and $\Pi_g$ (three bands); no coincidence.
\end{solution}

\begin{solution}{exo:b3:group-theory-applied:12}
$1b_2 \to 4a_1$: $B_2$, $y$, allowed, perpendicular to the plane. $3a_1 \to 4a_1$:
$A_1$, $z$, allowed, along the $C_2$ axis. $1b_2 \to 2b_1$: $A_2$, forbidden.
$1b_1 \to 4a_1$: $B_1$, $x$, allowed, in the plane perpendicular to the axis.
\end{solution}

\begin{solution}{pb:b3:group-theory-applied:1}
\textbf{1.} cis: the two CO on adjacent vertices; trans: opposite; fac: three
CO on one face; mer: three CO on a meridian (two trans, one between).
\textbf{2.} $C_{2v}$.
\textbf{3.} $D_{4h}$.
\textbf{4.} $C_{3v}$.
\textbf{5.} $C_{2v}$.
\textbf{6.} Only trans-\ce{ML4(CO)2}.
\textbf{7.} An operation that moves a CO onto another puts its vector off the
diagonal (contribution 0); one that leaves a CO in place maps its bond vector
onto itself (+1); none reverses a C--O vector.
\textbf{8.} \omterm{def:b3:point-groups:representation}{Characters} $(2, 0, 2, 0)$, with $\sigma_v(xz)$ the plane of both CO: $A_1 +
B_1$.
\textbf{9.} \omterm{def:b3:point-groups:representation}{Characters} 2 under $E$, $C_4$, $C_2$, $\sigma_v$, $\sigma_d$, 0 elsewhere:
$A_{1g} + A_{2u}$.
\textbf{10.} $(3, 0, 1)$: $A_1 + E$.
\textbf{11.} $(3, 1, 3, 1)$: $2A_1 + B_1$.
\textbf{12.} 2, 2, 3, 3: one dimension per CO.
\textbf{13.} cis 2, trans 1, fac 2, mer 3.
\textbf{14.} cis 2, trans 1, fac 2, mer 3.
\textbf{15.} trans: its IR band ($A_{2u}$) and Raman band ($A_{1g}$) are different
modes.
\textbf{16.} $A_{2u}$ changes sign under $i$: one CO lengthens while the other
shortens, out of phase.
\textbf{17.} $E$ is a degenerate pair: two modes of the same frequency give one
band, plus the $A_1$ band.
\textbf{18.} A trans-disubstituted, centrosymmetric isomer.
\textbf{19.} Three IR bands: mer (fac would give two).
\textbf{20.} Yes: mer has three Raman-active CO modes at the same frequencies,
being non-centrosymmetric; fac would show two.
\textbf{21.} $A_1$: the in-phase stretch is totally symmetric.
\textbf{22.} A fac/mer mixture would show up to five IR bands, with the fac
bands at different positions; three clean bands fit a single isomer. A
${}^{13}$C NMR spectrum would decide: two carbonyl signals in the ratio 2:1 for
mer, one signal for fac.
\textbf{23.} \textbf{The mer isomer has three IR-active CO stretches, the fac
isomer two}: the product is mer.
\end{solution}
